% Part: propositional-logic
% Chapter: propositional-logic
% Section: preliminaries

\documentclass[../../../include/open-logic-section]{subfiles}

\begin{document}

\olfileid[hi]{pl}{syn}{pre}

\olsection{प्रारंभिक तथ्य}

\begin{thm}[\emph{!!{formula}s पर आगमन का सिद्धांत}]
  \ollabel{thm:induction}  
  यदि कोई गुणधर्म~$P$ सभी परमाणु !!{formula}s के लिए सत्य है और ऐसा है कि
  \begin{enumerate}
    \tagitem{prvNot}{वह $\lnot !A$ के लिए सत्य हो, जब भी वह
      $!A$ के लिए सत्य हो;}{}
    \tagitem{prvAnd}{वह $(!A \land !B)$ के लिए सत्य हो,
      जब भी वह $!A$ और~$!B$ के लिए सत्य हो;}{}
    \tagitem{prvOr}{वह $(!A \lor !B)$ के लिए सत्य हो,
      जब भी वह $!A$ और~$!B$ के लिए सत्य हो;}{}
    \tagitem{prvIf}{वह $(!A \lif !B)$ के लिए सत्य हो,
      जब भी वह $!A$ और~$!B$ के लिए सत्य हो;}{}
    \tagitem{prvIff}{वह $(!A \liff !B)$ के लिए सत्य हो,
      जब भी वह $!A$ और~$!B$ के लिए सत्य हो;}{}
  \end{enumerate}
  तो $P$ सभी !!{formula}s के लिए सत्य है।
\end{thm}

\begin{proof}
  $S$ उन सभी !!{formula}s का संग्रह हो जिनमें गुणधर्म~$P$ है। स्पष्टतः
  $S \subseteq \Frm[L_0]$। $S$~\olref[fml]{defn:formulas} की सभी शर्तें
  पूरी करता है: उसमें सभी परमाणु !!{formula}s हैं और वह !!{operator}s के
  अंतर्गत संवृत है। $\Frm[L_0]$ ऐसा सबसे छोटा वर्ग है, इसलिए
  $\Frm[L_0] \subseteq S$। अतः $\Frm[L_0] = S$, और प्रत्येक सूत्र में
  गुणधर्म~$P$ है।
\end{proof}

\begin{prop}\ollabel{prop:balanced}
   प्रत्येक !!{formula}, जो~$\Frm[L_0]$ में हो, \emph{संतुलित} है: उसमें
   बाएँ और दाएँ कोष्ठकों की संख्या समान है।
\end{prop}

\begin{prob}
\olref[pl][syn][pre]{prop:balanced} सिद्ध करें।
\end{prob}

\begin{prop} \ollabel{prop:noinit}
!!a{formula} का कोई उचित आरंभिक खंड स्वयं !!a{formula} नहीं है।
\end{prop}

\begin{prob}
  \olref[pl][syn][pre]{prop:noinit} सिद्ध करें।
\end{prob}

\begin{prop}[अद्वितीय पठनीयता]
किसी भी !!{formula}~$!A$ का, जो $ \Frm[L_0]$ में हो, निम्नलिखित रूपों में
ठीक एक वाक्यविन्यास-विश्लेषण होता है:
\begin{enumerate}
\tagitem{prvFalse}{$\lfalse$।}{}

\tagitem{prvTrue}{$\ltrue$।}{}

\item $\Obj p_n$, जहाँ कोई $\Obj p_n \in  \PVar$ लिया गया हो।
  
\tagitem{prvNot}{$\lnot !B$, जहाँ कोई !!{formula}~$!B$ लिया गया हो।}{}

\tagitem{prvAnd}{$(!B \land !C)$, जहाँ !!{formula}s में से $!B$ और~$!C$ लिए गए हों।}{}

\tagitem{prvOr}{$(!B \lor !C)$, जहाँ !!{formula}s में से $!B$ और~$!C$ लिए गए हों।}{}

\tagitem{prvIf}{$(!B \lif !C)$, जहाँ !!{formula}s में से $!B$ और~$!C$ लिए गए हों।}{}

\tagitem{prvIff}{$(!B \liff !C)$, जहाँ !!{formula}s में से $!B$ और~$!C$ लिए गए हों।}{}
\end{enumerate}
इसके अतिरिक्त यह विश्लेषण \emph{अद्वितीय} है।
\end{prop}

\begin{proof}
$!A$ पर आगमन द्वारा। उदाहरण के लिए, मानें कि $!A$ के दो भिन्न विश्लेषण
$(!B \lif !C)$ और $(!B' \lif !C')$ हैं। तब $!B$ और $!B'$ समान होने चाहिए
(अन्यथा उनमें से एक दूसरे का उचित आरंभिक खंड होगा); अतः यदि $!A$ के दोनों
विश्लेषण भिन्न हों, तो इसका कारण यही होगा कि $!C$ और $!C'$ एक ही
प्रतीक-शृंखला के भिन्न विश्लेषण हैं, जो आगमनात्मक परिकल्पना से असंभव है।
\end{proof}


\begin{defn}[एकसमान प्रतिस्थापन]
यदि $!A$ और $!B$ !!{formula}s में से हों, और $\Obj p_i$ कोई !!{propositional
variable} हो, तो $\Subst{!A}{!B}{\Obj p_i}$ उस परिणाम को निरूपित करता है जो
$\Obj p_i$ की प्रत्येक उपस्थिति को $!B$ की एक उपस्थिति से $!A$ में बदलने पर
मिलता है; इसी प्रकार $\Obj p_1$, \dots,~$\Obj p_n$ का प्रतिस्थापन करने के लिए
!!{formula}s में से क्रमशः $!B_1$, \dots,~$!B_n$ लेकर युगपत परिणाम इस प्रकार
निरूपित किया जाता है:
$\SSubst{!A}{\subst{!B_1}{\Obj p_1},\dots,\subst{!B_n}{\Obj p_n}}$.
\end{defn}

\begin{prob} नीचे दिए पाँच !!{formula}s में से प्रत्येक के लिए निर्धारित करें
कि क्या वह !!{formula} किसी प्रतिस्थापन \( \Subst{!A}{!B}{\Obj p_i} \) के रूप
में व्यक्त किया जा सकता है, जहाँ \( !A \) क्रमशः (i) \( \Obj p_0 \);
(ii) \( ( \lnot \Obj p_0 \land \Obj
 p_1) \); और (iii) \( ( ( \lnot \Obj p_0 \lif \Obj p_1 ) \land \Obj
 p_2 ) \) है। प्रत्येक
स्थिति में उपयुक्त प्रतिस्थापन बताएँ।
  \begin{enumerate}
    \item \( \Obj p_1 \) 
    \item \( ( \lnot \Obj p_0 \land \Obj p_0 ) \)
    \item \( ( ( \Obj p_0 \lor \Obj p_1 ) \land \Obj p_2 )  \)
    \item \( \lnot ( ( \Obj p_0 \lif \Obj p_1 ) \land \Obj p_2 ) \)
    \item \( (( \lnot ( \Obj p_0 \lif \Obj p_1 ) \lif ( \Obj p_0 \lor \Obj p_1 )) \land \lnot ( \Obj p_0 \land \Obj p_1 )) \)
  \end{enumerate}
\end{prob}

\begin{prob}
$\Subst{!A}{!B}{p}$ की आगमन द्वारा गणितीय रूप से कठोर परिभाषा दें।
\end{prob}

\end{document}

